\chapter{定态近似方法}
\label{ch:approximations}

\section{非简并微扰}
将零场光谱已知的原子置于弱外场中，谱峰会随场强变化而移动。
以零场能量和本征态为起点，分别求出上下能级的一阶和二阶能移，便可预测频移对场强的斜率和曲率。
先考虑孤立的非简并能级。它与相邻能级的间隔决定了外场使本征态发生混合的程度。

定态微扰理论(time-independent perturbation theory)研究 $\op{H}(\lambda) = \op{H}_0 + \lambda \op{V}$。
其中 $\op{H}_0,\op{V}$ 是算符，$\op{V}$ 具有能量量纲，$\lambda$ 是调节扰动强度的无量纲实数。
取一个孤立的非简并本征态$\op{H}_0\ket{n} = E_n^{(0)}\ket{n}$，展开
\begin{subequations}
    \begin{align}
        E_n &= E_n^{(0)} + \lambda E_n^{(1)} + \lambda^2E_n^{(2)} + \cdots,\\
        \ket{n(\lambda)} &= \ket{n} + \lambda\ket{n^{(1)}} + \cdots.
    \end{align}
\end{subequations}

上标$(r)$表示$\lambda^r$的系数；
例如 $E_n^{(2)} = E_n''(0)/2$，其中撇表示对扰动强度 $\lambda$ 求导。
在未扰动本征基中，扰动算符 $\op{V}$ 的矩阵元为 $V_{mn} = \mel{m}{\op{V}}{n}$。每个 $V_{mn}$ 都是复数，对角元为实数。
将能量和态展开同时代入本征方程，
按$\lambda^0,\lambda^1,\lambda^2$收集系数：
\begin{subequations}\label{eq:pert-orders}
    \begin{alignat}{1}
        \op{H}_0\ket{n} &= E_n^{(0)}\ket{n},\\
        \op{H}_0\ket{n^{(1)}} + \op{V}\ket{n} &= E_n^{(0)}\ket{n^{(1)}} + E_n^{(1)}\ket{n},\\
        \op{H}_0\ket{n^{(2)}} + \op{V}\ket{n^{(1)}} &= E_n^{(0)}\ket{n^{(2)}}
            + E_n^{(1)}\ket{n^{(1)}} + E_n^{(2)}\ket{n}.
    \end{alignat}
\end{subequations}

中间归一化(intermediate normalization) $\braket{n}{n(\lambda)} = 1$固定了原态方向的系数。截断后的态还需按其范数归一化。
一级方程
\begin{equation}
    (\op{H}_0 - E_n^{(0)}\one)\ket{n^{(1)}} = (E_n^{(1)}\one - \op{V})\ket{n}
\end{equation}
左乘$\bra{n}$及$\bra{m}$（$m \ne n$），分别给出
\begin{subequations}
    \begin{align}
        E_n^{(1)} &= V_{nn},\\
        \ket{n^{(1)}} &= \sum_{m\ne n}\ket{m}\frac{V_{mn}}{E_n^{(0)} - E_m^{(0)}}.
    \end{align}
\end{subequations}

中间归一化使$\braket{n}{n^{(r)}} = 0$（$r \geq 1$）。
在式\eqref{eq:pert-orders}的二阶方程左乘$\bra{n}$，
两项$E_n^{(0)}\braket{n}{n^{(2)}}$抵消，
且 $E_n^{(1)}\braket{n}{n^{(1)}} = 0$。剩余的矩阵元可用一阶态修正展开，得到
\begin{equation}
    E_n^{(2)} = \mel{n}{\op V}{n^{(1)}}
        = \sum_{m\ne n}\frac{\abs{V_{mn}}^2}{E_n^{(0)} - E_m^{(0)}}.
    \label{eq:pertsecond}
\end{equation}
连续谱应加相应积分。对孤立的非简并基态，每个分母都为负，所以二阶修正非正。
一阶混合还会改变态的归一化。记 $s$ 为一阶修正态的范数平方，由正交性得到
\[
    s = \braket{n^{(1)}}{n^{(1)}}
      = \sum_{m\ne n}\biggabs{\frac{V_{mn}}{E_n^{(0)} - E_m^{(0)}}}^2 .
\]
中间归一化的态满足
\begin{equation}
    \braket{n(\lambda)}{n(\lambda)}
    = 1 + \lambda^2\braket{n^{(1)}}{n^{(1)}} + \bigo(\lambda^3).
\end{equation}
取范数平方的负二分之一次幂，归一化因子便是 $1 - \lambda^2s/2 + \bigo(\lambda^3)$。
这里写到$\lambda^2$的原态方向系数；
与$\ket{n}$正交的二阶分量仍由二阶投影方程确定。

对每个 $m \ne n$，把耦合 $\abs{\lambda V_{mn}}$ 与两态的能量差相比，就得到了一级混合振幅的大小。各振幅模平方之和衡量原态以外的总权重。
高阶级数可能是渐近展开(asymptotic expansion)。当一组能级趋于简并，应将它们一起纳入近似子空间。

\begin{example}{恒力作用下的谐振子}{qm-shifted-oscillator}
    在谐振子上加恒力 $F$，总势能是 $m\omega^2 \op{X}^2/2 - F\op{X}$。求基态的一级态修正、二阶能移和平均位置，并与经典平衡位置比较。

    \par\smallskip\textbf{解}\quad 把实际微扰 $\lambda \op{V}$ 取为 $-F\op{X}$。基态有 $\mel{0}{\op{X}}{0} = 0$，
    而 $\op{X}\ket{0} = x_0\ket{1}$，其中 $x_0 = \sqrt{\hbar/(2m\omega)}$。
    因此一阶能移为零；态修正与二阶能移各只剩一项：
    \begin{subequations}
        \begin{alignat}{1}
            \ket{0(F)} &= \ket{0} + \eta\ket{1} + \bigo(\eta^2),\\
            \eta &= \frac{-Fx_0}{E_0^{(0)} - E_1^{(0)}} = \frac{Fx_0}{\hbar\omega},\\
            \lambda^2 E_0^{(2)} &= -\frac{F^2x_0^2}{\hbar\omega}
                = -\frac{F^2}{2m\omega^2}.
        \end{alignat}
    \end{subequations}
    用一阶态计算 $\ave{\op{X}}$，两个交叉项之和为 $2\eta x_0 = F/(m\omega^2)$，
    恰好给出平衡位置的移动。配平方还能独立检验能量：
    \begin{equation}
        \frac{m\omega^2 \op{X}^2}{2} - F\op{X}
            = \frac{m\omega^2}{2}\biggp{\op{X} - \frac{F}{m\omega^2}}^2
                  - \frac{F^2}{2m\omega^2}.
    \end{equation}
    势阱平移使每一级的能量都减少同一个常数。能级间距保持 $\hbar\omega$，
    但位置测量已能在 $F$ 的一阶读出变化。态展开要求 $\abs{\eta} \ll 1$。
\end{example}

\section{简并、近简并与避免交叉}
上一节的态修正含有两个能级的能量差作分母。若未扰动能量相等，即使把外场调得任意小，这个分母仍为零。
原本重合的谱峰在弱场下如何分裂，可由各条分支在零场处的斜率刻画。求这些斜率首先需要找出外场在简并子空间中选定的线性组合。

设 $E_0$ 是 $\op H_0$ 的简并本征值。记 $\op\Pi$ 为投影到相应简并子空间(degenerate subspace)的正交投影算符，$\op\Pi_\perp = \one - \op\Pi$ 则投影到它的正交补。
设 $\ket{\phi}$ 是简并子空间内待选的零级态。以 $\op\Pi$ 左乘一级方程时，$\op\Pi(\op H_0 - E_0\one) = 0$ 消去了未知的一阶态，留下
\begin{equation}
    \op\Pi\op V\op\Pi\ket{\phi} = E^{(1)}\ket{\phi}.
\end{equation}
因此先在简并子空间内对角化 $\op\Pi\op V\op\Pi$，这个较小矩阵的本征矢给出零级态，本征值给出一级能移。

将精确态 $\ket{\psi}$ 分解为简并子空间内的分量 $\ket{\psi_\parallel} = \op\Pi\ket{\psi}$ 和正交补内的分量 $\ket{\psi_\perp} = \op\Pi_\perp\ket{\psi}$。
由于$\op\Pi + \op\Pi_\perp = \one$、$\op\Pi\op\Pi_\perp = 0$，精确本征方程分成
\begin{subequations}
    \begin{alignat}{1}
        \op\Pi\op{H}\op\Pi\ket{\psi_\parallel} + \op\Pi\op{H}\op\Pi_\perp\ket{\psi_\perp} &= E\ket{\psi_\parallel},\\
        \op\Pi_\perp\op{H}\op\Pi\ket{\psi_\parallel} + \op\Pi_\perp\op{H}\op\Pi_\perp\ket{\psi_\perp} &= E\ket{\psi_\perp}.
    \end{alignat}
\end{subequations}
在正交补上，单位算符记为 $\one_\perp$。若 $E\one_\perp - \op\Pi_\perp\op{H}\op\Pi_\perp$ 在该空间上可逆，第二式给出
\begin{equation}
    \ket{\psi_\perp} = (E\one_\perp - \op\Pi_\perp\op{H}\op\Pi_\perp)^{-1}\op\Pi_\perp\op{H}\op\Pi\ket{\psi_\parallel}.
\end{equation}
代回第一式，得到作用在保留子空间中的有效 Hamilton 算符
\begin{equation}
    \op{H}_\eff(E) = \op\Pi\op{H}\op\Pi + \op\Pi\op{H}\op\Pi_\perp(E\one_\perp - \op\Pi_\perp\op{H}\op\Pi_\perp)^{-1}\op\Pi_\perp\op{H}\op\Pi.
\end{equation}
对隔离简并簇展开到二阶，
\begin{equation}
    \op{H}_\eff = E_0\op\Pi + \lambda \op\Pi\op{V}\op\Pi\\
     +\lambda^2\op\Pi\op{V}\op\Pi_\perp(E_0\one_\perp - \op\Pi_\perp\op{H}_0\op\Pi_\perp)^{-1}
        \op\Pi_\perp\op{V}\op\Pi + \bigo(\lambda^3).
\end{equation}
一级仍简并时还须对角化二级矩阵。
当两态近简并时，保留整个两态子空间，在其基底中直接对角化
\begin{subequations}
    \begin{align}
        \op{H}_2 &= \begin{pmatrix}E_a&g\\g\conj&E_b\end{pmatrix},\\
        E_\pm &= \frac{E_a + E_b}{2} \pm \sqrt{\frac{(E_a - E_b)^2}{4} + \abs{g}^2}.
    \end{align}
\end{subequations}
非零耦合产生避免交叉(avoided crossing)：在$E_a = E_b$处劈裂为$2\abs{g}$（图\ref{fig:avoided-crossing}）；若对称性强制$g = 0$，两支可以相交。
远离交叉展开根号，正好恢复式\eqref{eq:pertsecond}，这是一个直接的内部检验。
\begin{figure}[htbp]
    \centering\includegraphics[width=.60\linewidth]{figures/avoided_crossing.pdf}
    \caption{点线为未耦合能级，实线为耦合后的本征值。纵轴以 $\bar E = (E_a + E_b)/2$ 为零点、以 $\abs{g}$ 为单位；在两条点线的交点处量取实线的竖直间距，即得最小能隙。}
    \label{fig:avoided-crossing}
\end{figure}

\begin{example}{远离交叉的两能级}{qm-avoided-crossing}
    两个裸能级 $E_a < E_b$ 的耦合矩阵元为 $g = \mel{a}{\op H}{b}$。当 $\abs{g} \ll \Delta = E_b - E_a$ 时，求低能支的能移和另一裸态混入的振幅，各保留最低非零阶。

    \par\smallskip\textbf{解}\quad 将两能级的精确本征值按 $\abs{g}/\Delta$ 展开，得到
    \begin{equation}
        \begin{alignedat}{1}
            E_- &= \frac{E_a + E_b}{2}
                - \frac{\Delta}{2}\sqrt{1 + \frac{4\abs{g}^2}{\Delta^2}}\\
            &= E_a - \frac{\abs{g}^2}{\Delta}
                + \bigo\Bigp{\frac{\abs{g}^4}{\Delta^3}}.
        \end{alignedat}
    \end{equation}
    同时本征方程的第二行给出
    $\braket{b}{-}/\braket{a}{-} \simeq -g\conj/\Delta$。
    能移是二阶，态混合已经是一阶。
\end{example}

\begin{example}{三能级系统的简并劈裂}{qm-degenerate-dark-combination}
    三个正交归一态 $\ket a,\ket b,\ket e$ 的未扰动能量分别为 $0,0,\Delta$，其中 $\Delta > 0$。
    微扰仅有 $\mel{e}{\op V}{a} = g$、$\mel{e}{\op V}{b} = 2g$ 及其共轭矩阵元非零，$g$ 为实数，$\abs g \ll \Delta$。
    低能双态从哪一阶开始劈裂？有没有任意 $g$ 下都精确保持零能量的态？

    \par\smallskip\textbf{解}\quad 简并子空间内的微扰矩阵全为零，一级无法选定基底。考虑任意组合
    $c_a\ket a + c_b\ket b$，它通向 $\ket e$ 的振幅正比于 $c_a + 2c_b$。
    因此 $\ket D = (2\ket a - \ket b)/\sqrt5$ 被 $\op V$ 消去，总 Hamilton 算符作用在它上面也给出零。
    正交组合 $\ket B = (\ket a + 2\ket b)/\sqrt5$ 与 $\ket e$ 的耦合为 $\sqrt5\,g$。
    在 $\ket B,\ket e$ 子空间对角化，低能一支为
    \[
        E_- = \frac{\Delta - \sqrt{\Delta^2 + 20g^2}}{2}
        = -\frac{5g^2}{\Delta} + \bigo\bigp{g^4/\Delta^3}.
    \]
    两支因此在二阶分开。若分别对 $\ket a,\ket b$ 套非简并二阶公式，只会算到二阶矩阵的两个对角元，
    还会漏掉经同一中间态产生的非对角元 $-2g^2/\Delta$。
\end{example}

\section{Stark效应}
\label{sec:stark}
静电场中的原子提供了区分简并与非简并微扰的实验例子。在弱场区，有些能级的位移正比于场强的平方，
另一些原本重合的能级在很小的场中就线性分裂。这些静电场引起的能移称为 Stark 效应。
判断曲线从哪一阶开始，要同时知道未扰动态的宇称，以及是否有别的态与它简并。

沿$z$的静电场$\mathcal E$给出$\op{V} = -\mathcal E \op{d}_z$。
若零场态非简并且有确定宇称，因为$\op{d}_z$是奇算符，一阶能移为零。
二阶则为
\begin{subequations}
    \begin{align}
        \Delta E_n &= -\frac12\alpha_n\mathcal E^2,\\
        \alpha_n &= 2\sum_{m\ne n}\frac{\abs{\mel{m}{\op{d}_z}{n}}^2}{E_m^{(0)} - E_n^{(0)}}.
    \end{align}
\end{subequations}
这里$\alpha_n$是该方向的静极化率(static polarizability)，SI单位为$\unit{\coulomb\metre\squared\per\volt}$；基态的$\alpha_0 > 0$。
激发态可以耦合到更高或更低的能级，两类中间态对二阶能移的贡献符号相反。

二阶能量中的$1/2$可以通过诱导偶极核对。
由Hellmann--Feynman关系，
\begin{equation}
    \ave{\op{d}_z}_n = -\pv{E_n}{\mathcal E}
        = \alpha_n\mathcal E + \bigo(\mathcal E^2).
\end{equation}
逐渐开启电场时，偶极本身从零增长到 $\alpha_n\mathcal E$。
把每次微小的能量改变量 $-\ave{\op{d}_z}_n\d\mathcal E$ 从零场累加到终场，得到二阶能移
\[
    \Delta E_n^{(2)} = -\int_0^\mathcal E\alpha_n\mathcal E'\d\mathcal E'
    = -\frac12\alpha_n\mathcal E^2 .
\]
把最终的诱导偶极矩乘以 $-\mathcal E$ 将得到实际二阶能移的两倍。

若一对相反宇称态近简并，且$\mel{a}{\op{d}_z}{b} = d$，
其有效矩阵的非对角元为$-\mathcal E d$。
在精确简并极限，一阶能移为$\pm\abs{\mathcal E d}$，本征态成为带永久偶极(permanent dipole)的叠加。
氢原子同一$n$壳层的相反宇称态提供了线性Stark效应的典型场景；
真实氢中精细结构与Lamb位移已先移除部分简并，应比较场致耦合与这些能量差再选择线性或二次区间。
无限均匀静电场还允许电离，严格束缚能级变成共振(resonance)；微弱电场下的Stark展开描述寿命足够长的亚稳能级。

\begin{example}{相反宇称双态的 Stark 劈裂}{qm-stark-crossover}
    取一对相反宇称态，零场间距为$\Delta > 0$，
    并选择相位使偶极矩阵元$d$为实数。
    以两态平均能量$\avg E$为零场中心，矩阵为
    \begin{equation}
        \op{H} = \avg E\,\one +
            \begin{pmatrix}
                -\Delta/2&-d\mathcal E\\
                -d\mathcal E&\Delta/2
            \end{pmatrix}.
    \end{equation}
    求两个能级的能移，并比较 $\abs{d\mathcal E}$ 远小于和远大于 $\Delta$ 时的场强依赖。

    \par\smallskip\textbf{解}\quad 解行列式 $\det(\op{H} - E\one) = 0$，得到
    \begin{equation}
        E_\pm = \avg E \pm
            \sqrt{\frac{\Delta^2}{4} + d^2\mathcal E^2}.
    \end{equation}
    当$\abs{d\mathcal E} \ll \Delta$，低能支移动
    $-d^2\mathcal E^2/\Delta$，对应$\alpha = 2d^2/\Delta$。
    当$\Delta \ll \abs{d\mathcal E}$且仍未混入其他态时，
    劈裂近似为$2\abs{d\mathcal E}$。
    同一模型在不同场强区间给出不同的表观场依赖，
    因而“线性”或“二次”Stark效应应连同所比较的能隙说明。
\end{example}

\section{变分法}

Stark 能移的微扰展开使用已知能级和偶极矩阵元。若只知道势能函数而难以求得精确本征态，还能怎样估计基态能量？
例如，可选取一族宽度连续可调的归一化波函数，计算能量期望随宽度的变化，求出其中的最小值。
由于所有态的能量期望均受能谱下界约束，这个最小值给出了基态能量的上界。

将归一化试探态 $\ket\phi$ 的能量测量分布记为 $\mu_\phi$，即 $\mu_\phi(B)$ 是能量读数落入区间 $B$ 的概率。
\begin{theorem}[unbreakable]{变分上界}{qm-variational}
    设自伴 Hamilton 算符 $\op H$ 的能谱下确界为 $E_0$。对能量均值有限的任意归一化试探态，
    \begin{equation}
        E_{\trial} = \int E\,\d\mu_\phi(E) \geq E_0 .
    \end{equation}
    当$E_0$为孤立基态能量时，等号要求试探态(trial state)完全位于基态子空间。
\end{theorem}
\begin{proof}
    能量分布的总概率为 $1$，且只在 $E \ge E_0$ 处有权重，所以
    \begin{equation}
        E_{\trial} - E_0
        = \int(E - E_0)\d\mu_\phi(E) \geq 0.
    \end{equation}
    等号要求在$E > E_0$的部分没有概率权重。
\end{proof}
若 $\op H\ket\phi$ 也是态空间中的向量，就可直接用矩阵元 $\mel{\phi}{\op H}{\phi}$ 计算这个均值。

变分法(variational method)给出所选试探族中的最低能量，它是基态能量的上界。上界的精度取决于试探族逼近基态的能力。
激发态的上界则由极小极大原理(min--max principle)给出；
直接变分时，应将试探态限制在与真实低能态正交的子空间。

以$\op{H} = \op{P}^2/(2m) + m\omega^2\op{X}^2/2 + \lambda \op{X}^4$、$\lambda > 0$为例。
取频率为 $\Omega > 0$ 的辅助谐振子真空 $\ket{0(\Omega)}$ 作为试探态，升降代数给出
\begin{equation}
    E_{\trial}(\Omega) = \frac{\hbar\Omega}{4}
    + \frac{\hbar\omega^2}{4\Omega}
    + \frac{3\lambda\hbar^2}{4m^2\Omega^2}.
\end{equation}

式中的三项分别使用
\begin{subequations}
    \begin{align}
        \ave{\op{P}^2}_\Omega &= \frac{m\hbar\Omega}{2},\\
        \ave{\op{X}^2}_\Omega &= \frac{\hbar}{2m\Omega},\\
        \ave{\op{X}^4}_\Omega &= \frac{3\hbar^2}{4m^2\Omega^2}.
    \end{align}
\end{subequations}
最后一项由\exmref{exm:qm-x-fourth}取$n = 0$得到。
对$\Omega$求导，
\[
    \dv{E_{\trial}}\Omega
    = \frac{\hbar}{4} - \frac{\hbar\omega^2}{4\Omega^2}
    - \frac{3\lambda\hbar^2}{2m^2\Omega^3}.
\]

驻值条件是
$\Omega^3 - \omega^2\Omega - 6\lambda\hbar/m^2 = 0$。
$\lambda \to 0$时$\Omega \to \omega$，并恢复一级能移
$3\lambda\hbar^2/(4m^2\omega^2)$。
此例中的$\lambda$带量纲，和本章开头的无量纲记账参数不同；小量为$\lambda\hbar/(m^2\omega^3)$。

\begin{example}{一维弱吸引势中束缚态的存在性}{qm-exam-weak-binding}
    粒子在整条直线上运动，$V(x) = -\lambda U(x)$，其中 $\lambda > 0$，$U$ 非负、有界且只在有限区间内非零，积分 $S = \int_\RR U(x)\d x > 0$。
    已知散射连续谱从零开始。证明无论吸引强度 $\lambda$ 多小，体系都至少有一个负能束缚态。

    \par\smallskip\textbf{解}\quad 由变分上界，只要找到一个能量均值为负的归一化态即可。取宽度可变的 Gauss 波包
    \[
        \psi_R(x) = (\pi R^2)^{-1/4}\exp[-x^2/(2R^2)] .
    \]
    当 $R$ 远大于势的范围，波包在势区内的概率密度近似为 $1/(\sqrt\pi R)$。由 $\psi_R' = -x\psi_R/R^2$ 和 $\ave{X^2} = R^2/2$ 得
    \begin{align*}
        \ave{T}_R &= \frac{\hbar^2}{4mR^2},\\
        \ave{V}_R &= -\frac{\lambda}{\sqrt\pi R}\int_\RR U(x)\e^{-x^2/R^2}\d x
                   = -\frac{\lambda S}{\sqrt\pi R} + o(R^{-1}) .
    \end{align*}
    动能以 $1/R^2$ 下降，吸引势能的绝对值只以 $1/R$ 下降。对每个固定的 $\lambda > 0$，取足够大的 $R$ 就使总能量均值低于连续谱阈值，于是阱中至少有一个束缚态。
\end{example}

\subsection{氦原子的基态与电子屏蔽}
\label{ss:qm-helium-ground}
氦离子 $\mathrm{He}^+$ 的一个电子束缚在电荷为 $+2e$ 的核上。加入第二个电子后，
两个电子都受核吸引，又彼此排斥。电子云要扩展到多大，才能使总能量最低？
取固定的点核，忽略相对论和辐射修正，记两电子到核的距离为 $r_1,r_2$，彼此的距离为 $r_{12} = \abs{\bm r_1 - \bm r_2}$。
用 Bohr 半径 $a_0$ 和 Hartree 能量 $E_{\mathrm h} = \hbar^2/(m_\electron a_0^2) = e^2/(4\pi\varepsilon_0 a_0)$ 表示位置表象中的 Hamilton 算符，有
\begin{equation}\label{eq:qm-helium-scaled-hamiltonian}
    \frac{\op H}{E_{\mathrm h}}
    = \sum_{i=1}^2\biggp{-\frac{a_0^2}{2}\nabla_i^2 - \frac{2a_0}{r_i}}
        + \frac{a_0}{r_{12}} .
\end{equation}
最后一项把两个位置变量耦合起来。先让两个电子共用一个氢样 $1\mathrm s$ 轨道，把轨道的衰减长度 $a_0/\zeta$ 作为可调参数\cite{mithelium}。
归一化的空间试探态为
\begin{equation}\label{eq:qm-helium-product-trial}
    \Phi_\zeta(\bm r_1,\bm r_2) = \phi_\zeta(r_1)\phi_\zeta(r_2),\qquad
    \phi_\zeta(r) = \sqrt{\frac{\zeta^3}{\pi a_0^3}}\,\e^{-\zeta r/a_0},\qquad \zeta > 0 .
\end{equation}
这个 $1\mathrm s^2$ 组态的空间因子在交换下对称，乘以反对称自旋单态
$(\ket{\up,\dn} - \ket{\dn,\up})/\sqrt2$，总态就满足电子的交换要求（\ssecref{ss:qm-helium-exchange}）。
算符 $\op H$ 与自旋无关，求能量只需对空间部分积分。

对单个电子，$\abs{\grad\phi_\zeta}^2 = (\zeta/a_0)^2\abs{\phi_\zeta}^2$ 给出动能 $\zeta^2 E_{\mathrm h}/2$，
径向积分给出 $\ave{a_0/r} = \zeta$。求斥力能时，固定 $r_1,r_2$，对两径向方向的夹角 $\theta$ 作平均。
令 $u = \cos\theta$，则
\[
    \frac12\int_{-1}^1\frac{\d u}{\sqrt{r_1^2 + r_2^2 - 2r_1r_2u}}
    = \frac{r_1 + r_2 - \abs{r_1 - r_2}}{2r_1r_2}
    = \frac{1}{\max(r_1,r_2)} .
\]
引入 $x = 2\zeta r_1/a_0$、$y = 2\zeta r_2/a_0$，每个电子的径向概率密度都是 $p(x) = x^2\e^{-x}/2$。
两个积分区 $x > y$ 与 $y > x$ 的贡献相等，因而
\[
    \biggave{\frac{a_0}{r_{12}}}
    = \zeta\int_0^\infty y^2\e^{-y}\d y\int_y^\infty x\e^{-x}\d x
    = \zeta\int_0^\infty (y^2 + y^3)\e^{-2y}\d y
    = \frac{5\zeta}{8}.
\]
合并两电子的动能、核吸引能和电子斥力能，得到
\begin{equation}\label{eq:qm-helium-variational-energy}
    \frac{E_{\trial}(\zeta)}{E_{\mathrm h}}
       = \zeta^2 - 4\zeta + \frac{5\zeta}{8}
       = \biggp{\zeta - \frac{27}{16}}^2 - \frac{729}{256}.
\end{equation}
最低点 $\zeta = 27/16 < 2$ 对应的电子云比 $\mathrm{He}^+$ 的 $1\mathrm s$ 轨道更宽。$\zeta$ 是描述平均屏蔽的有效核电荷数，
Hamilton 算符中的实际核电荷始终是 $+2e$。乘积试探态可以调整电子云的大小，两个电子的位置概率却仍然相互独立。
变分结果还能检验两个电子能否一起被束缚。一个电子以零动能远离原子、另一个留在 $\mathrm{He}^+$ 基态时，总能量为 $-2E_{\mathrm h}$。
变分上界给出 $E_0 \le -729E_{\mathrm h}/256$，已经低于这个第一电离阈值。
所以在这个模型中，从基态移走一个电子至少需要 $(217/256)E_{\mathrm h} \simeq \qty{23.07}{eV}$ 的能量。
中性氦的第一电离能约为 $\qty{24.59}{eV}$\cite{nisthelium}。

\section{Hellmann--Feynman关系与残差检验}
\label{sec:hellmann}
变分计算给出能量上界。要与谱线随外场变化的斜率比较，还需确定能量的参数导数，并判断近似态满足本征方程的精度。
对本征方程求导可得到谱线斜率；把候选能量与候选态代回本征方程，则可用残差向量的范数度量偏差。

\begin{theorem}{Hellmann--Feynman定理}{qm-hellmann-feynman}
    设 $\op{H}(\eta)$ 是依赖实参数 $\eta$ 的自伴 Hamilton 算符，
    $\ket{n(\eta)}$ 为选定本征值分支上可微的归一化本征态，且本征方程的各项可对 $\eta$ 求导。则
    \begin{equation}
        \dv{E_n}{\eta} = \mel{n}{\dv{\op H}{\eta}}{n}.
    \end{equation}
\end{theorem}
\begin{proof}
用撇表示对 $\eta$ 求导，特别地 $\ket{n'} = \d\ket{n}/\nd\eta$。对 $\op{H}\ket{n} = E_n\ket{n}$ 求导，再左乘 $\bra{n}$，
\begin{alignat*}{1}
    \op H'\ket n + \op{H}\ket{n'} &= E_n'\ket n + E_n\ket{n'},\\
    \mel{n}{\op H'}{n} + E_n\braket{n}{n'} &= E_n' + E_n\braket{n}{n'} .
\end{alignat*}
第二行利用$\op{H}$自伴、$E_n$为实数及$\braket{n}{n} = 1$。
态随参数的变化项因而相消。
\end{proof}

简并点须先选择使参数导数在简并空间内对角的分支。
把未经驻值优化的任意试探态代入则通常有额外的态导数项。

若要检查候选态，将本征方程的两侧作差。
\begin{definition}{本征方程的残差}{qm-eigen-residual}
    若归一化试探态 $\ket{\phi}$ 的能量二阶矩有限，则 $\op H\ket\phi$ 是态空间中的向量。将试探态和候选能量 $E_{\trial}$ 代入本征方程，两侧之差称为残差向量(residual)。其范数
    \begin{equation}
        r = \bignorm{(\op H - E_{\trial}\one)\ket{\phi}}
    \end{equation}
    具有能量量纲。
\end{definition}
沿用试探态的能量分布 $\mu_\phi$，残差平方是能量偏差平方的平均值，
\begin{equation}\label{eq:qm-spectral-residual}
    r^2 = \int (E - E_{\trial})^2\d\mu_\phi(E).
\end{equation}
各谱值离 $E_{\trial}$ 的距离都至少等于它到谱的距离，因此这个最短距离不超过 $r$。
要从这个邻域内的谱识别出目标能级，还需要相邻能隙的信息。
若基态非简并、与其他谱的间隔至少为 $\Delta > 0$，且试探态的能量期望为$E_0 + \epsilon$，则
$1 - \abs{\braket{0}{\phi}}^2 \leq \epsilon/\Delta$。

因为在基态正交补上$E - E_0 \geq \Delta$，
\begin{equation}
    \epsilon
    = \int(E - E_0)\d\mu_\phi(E)
    \geq \Delta\bigp{1 - \abs{\braket{0}{\phi}}^2}.
\end{equation}
这把能量上界与态的重叠联系起来；
$\epsilon/\Delta$越小，试探态落在激发态子空间的概率越小。

\begin{remark}[截断基底中的残差]
    设 $\op\Pi$ 投影到计算所保留的基矢张成的子空间，试探态满足 $\op\Pi\ket{\phi} = \ket{\phi}$。
    即使已把 $\op\Pi\op{H}\op\Pi$ 精确对角化，$\op{H}\ket{\phi}$ 仍可能有落在截断空间外的分量
    $(\one - \op\Pi)\op{H}\ket{\phi}$。核验能级时，应把这部分也计入残差。
\end{remark}

\begin{example}{微扰下的能量与态混合}{qm-energy-and-state-orders}
    一个孤立能级的一阶能移为零，微扰还会怎样改变态和能量？若邻近还有相反宇称的能级，如何判断能否沿用同一展开？

    \par\smallskip\textbf{解}\quad 态混合通常已经是一阶，能量可从二阶改变。
    若微扰与相反宇称态之间的耦合可与原能隙相比，应将两态一起对角化，此时还可能产生近线性劈裂。
\end{example}
\reading{定态微扰与简并处理核对MIT 8.06讲义1\cite{mitpert}。
    下一章把外场变成随时间变化的探针，连接能级和实际吸收发射率。}
